\section{Глава~\ref{sec:debugging}. Отладка машины}

Что слышит электрод, маломерность активности, работа интерфейсов на жёстких решётках, совместное обучение мозга и декодера и зрительные точки от стимуляции измерены в рецензируемых опытах; оценки потоков данных --- арифметика главы; о вживлённом людям устройстве Neuralink, насколько удалось проверить, публиковала данные в основном сама компания.

{\footnotesize
\setlength{\tabcolsep}{3pt}\renewcommand{\arraystretch}{1.15}
\begin{longtable}{>{\raggedright}p{0.5cm}>{\raggedright}p{4.0cm}>{\raggedright}p{5.6cm}>{\raggedright}p{2.3cm}>{\raggedright\arraybackslash}p{1.7cm}}
\toprule
№ & Утверждение & Опыт и что измерено & Источник & Статус \\
\midrule\endfirsthead
\toprule
№ & Утверждение & Опыт и что измерено & Источник & Статус \\
\midrule\endhead
1 & Электрод слышит импульсы нейронов в радиусе порядка сотни микрометров, обычно от одного до нескольких; их различают по форме & Крыса, поле CA1, одновременная внутриклеточная и внеклеточная запись: форма внеклеточного импульса повторяет ширину и амплитуду внутриклеточного. Оценка: тетрод мог бы различать $60$--$100$ нейронов, а на деле выделяют около шести. & \cite{Henze2000extra} & измерено \\
2 & В мозге $8.6\cdot10^{10}$ нейронов; щуп слышит примерно одну стомиллионную & Подсчёт ядер во взвеси растворённой ткани: $86\cdot10^9$ нейронов в мозге человека. Доля --- арифметика главы. & \cite{Azevedo2009} & измерено \\
3 & Активность сотен нейронов двигательной коры укладывается в десяток-другой общих переменных & Обзор записей в двигательной коре обезьян: активность популяции описывается немногими скрытыми переменными, общими для многих нейронов. & \cite{Gallego2017} & измерено \\
4 & Шум соседей согласован, и сотня нейронов несёт почти столько же, сколько тысяча (утверждение~\ref{prop:pooling}) & Зрительная кора обезьяны: коэффициент корреляции шума соседних нейронов около $0.1$, и выигрыш от усреднения насыщается при сотне нейронов. Теория ограничивающих информацию корреляций. & \cite{Zohary1994, MorenoBote2014} & измерено \\
5 & Сырой поток $2\cdot10^8$~бит/с против $8\cdot10^4$~бит/с в импульсах~\eqref{eq:probedata} & Арифметика главы по ёмкости потока импульсов~\eqref{eq:spikeentropy}. Параметры оцифровки реальны: кристалл Neuralink --- $19.3$~кГц, $10$~бит на канал. & вывод в главе; \cite{Rieke1997, Musk2019} & теория \\
6 & Импульсы выделяют на кристалле у электродов и передают наружу только события; корпус закрыт кожей, связь и питание без проводов & Система, описанная в статье компании, передаёт полный сырой поток по кабелю USB-C, импульсы выделяет программа в реальном времени вне кристалла; сжатие на борту, питание и передача без проводов названы там планом для клинических устройств. Для вживлённого людям устройства --- только сообщения компании. & \cite{Musk2019}; отчёты Neuralink, рецензируемой статьи нет & измерено (отчёт компании) \\
7 & До $3072$ электродов на $96$ нитях по $32$; нити вставляет робот, обходя сосуды & Статья компании: $3072$ электрода на $96$ нитях по $32$; робот вставляет $6$ нитей ($192$ электрода) в минуту; до $70\,\%$ электродов у крыс с долговременно вживлённой решёткой слышат импульсы. & \cite{Musk2019} & измерено (статья компании) \\
8 & У людей около тысячи каналов на $64$ нитях; часть нитей сместилась; курсор порядка $10$~бит/с & Только сообщения компании. Поиск в PubMed не нашёл рецензируемой статьи с результатами на людях; это согласуется с оговоркой в тексте главы. & отчёты Neuralink; рецензируемой статьи нет & измерено (отчёт компании) \\
9 & Интерфейс на решётке из сотни электродов управляет курсором & Человек с параличом, $96$ электродов в первичной двигательной коре: намерение двигать рукой меняет импульсы через три года после травмы; декодер даёт «нейронный курсор» (почта, телевизор), управление протезом кисти и роботом-рукой. & \cite{Hochberg2006} & измерено \\
10 & Письмо «мысленной рукой»: около $90$ знаков в минуту, меньше $8$~бит/с & Человек с параличом кисти, попытки писать от руки, рекуррентный декодер: $90$ знаков в минуту при точности $94.1\,\%$ в реальном времени, выше $99\,\%$ после автокоррекции. Оценка в битах --- арифметика главы. & \cite{Willett2021} & измерено \\
11 & Попытки говорить переводятся в текст со скоростью около $60$ слов в минуту & Человек, потерявший внятную речь, решётки в коре: $62$ слова в минуту; $9.1\,\%$ ошибок в словах при словаре из $50$ слов, $23.8\,\%$ при словаре из $125\,000$. & \cite{Willett2023} & измерено \\
12 & Интерфейс вынимает малую долю ёмкости: у одного нейрона десятки бит в секунду, у сотни --- тысячи (утверждение~\ref{prop:probe}) & Верхняя граница~\eqref{eq:spikeentropy} --- теория; сравнение со строками~8, 10 --- вывод главы. Что узкое место --- маломерность и незнание кода, а не провод, прямым опытом не проверено. & \cite{Rieke1997} & теория \\
13 & Мозг и декодер учатся друг у друга & Обезьяны управляют курсором; декодер подстраивается в замкнутой петле, а нейроны одновременно переучиваются: точность растёт, навык сохраняется и меньше страдает от движений собственной руки. Совет развести скорости обучения --- инженерное правило, отдельного опыта в главе нет. & \cite{Orsborn2014coad} & измерено \\
14 & Ток через электрод возбуждает окружающие нейроны и провода без разбора типа & Кора мыши и кошки, двухфотонная запись кальция: даже слабый ток возбуждает редкие нейроны, разбросанные на расстояние до миллиметров, через прямое возбуждение аксонов в объёме диаметром в десятки микрометров. То есть возбуждение неизбирательно, но не сплошное. & \cite{Histed2009stim} & измерено \\
15 & Стимуляция зрительной коры зажигает точки; до $16$ точек сразу позволяют узнать некоторые буквы и границы предметов & Слепой человек, $96$ электродов в зрительной коре на $6$ месяцев: порог одного электрода $66.8\pm36.5$~мкА; одновременная стимуляция групп (до $16$ электродов --- предел стимулятора) снижает порог и даёт различимые узоры, некоторые буквы и границы предметов узнаются. & \cite{Fernandez2021} & измерено \\
16 & Второе направление Neuralink --- устройство, подающее сигналы в зрительную кору & Только сообщения компании о разработке; рецензируемых данных о его работе не найдено. & отчёты Neuralink & гипотеза \\
17 & Концентрации широковещательных медиаторов можно мерить химическим датчиком в ткани & Срезы мозга крысы, быстрая циклическая вольтамперометрия на углеродном волокне: измерены выброс и обратный захват серотонина; вещество выходит за пределы синапса. & \cite{Bunin1998} & измерено \\
\bottomrule
\end{longtable}}
